\chapter{Einf"uhrung}
Dieses Kapitel will in die Thematik des \www und die der Suchmaschinen einf"uhren. Das Themengebiet wird "uberblickhaft umrissen, die wichtigsten Begriffe werden erl"autert und die vorhandenen Probleme werden aufgezeigt.


\section{Konzept des \www}
Das \www wurde eigentlich f"ur die Verbreitung von Dokumenten am CERN\footnote{http://www.cern.ch/} Ende 1990 von Tim Berners-Lee\footnote{http://www.w3.org/pub/WWW/People/Berners-Lee/} entwickelt. Es war daf"ur vorgesehen, Dokumente der dort arbeitenden Wissenschaftler untereinander zug"anglich zu machen.

Das eigentlich Neue am \www war, da"s zum einen Dokumente mit einheitlichen Namen netzwerkweit zug"anglich wurden und zum anderen, da"s man Referenzen auf andere Dokumente mit aktiven Verweisen sog. {\it Links} erzeugen konnte. Es wurde somit sehr einfach, anstelle der Literaturhinweise, nun die referenzierten Dokumente einfach als aktiven Verweis in das eigene Dokument aufzunehmen.

Solche {\em Hypertext} genannten Konzepte zur Verkn"upfung von Dokumenten existieren schon l"anger ( z.B. HyperCard auf Macintosh Systemen ). Sowohl die Einfachheit der Bedienung, als auch die freie Verf"ugbarkeit der Anzeigeprogramme auf verschiedenen Rechnerplattformen im Internet f"uhrten dann zur weiten Verbreitung des \www.


\section{Konzept der Suchmaschinen}
Sowohl eine Informationsaufbereitung als auch eine "Uberwachung der im Dokument befindlichen Verweise fand von Hause aus im \www nicht statt, weshalb bald Suchmaschinen dieses Informationsloch zu stopfen suchten.

Suchmaschinen, beginnend an einem oder mehreren Startdokumenten, indexieren den in dem Dokument vorhandenen Text. Um neue Dokumente holen zu k"onnen, werten sie die in dem Dokument befindlichen Verweise aus. Nach bestimmten Regeln werden dann die Dokumente, auf welche verwiesen wurde, geholt, z.B. nur Seiten aus Deutschland. 

Suchmaschinen indexieren so nach und nach ganze Myriaden von Dokumenten. {\em AltaVista}\footnote{http://www.altavista.digital.com/} indexiert derzeit rund 50 Mio. Dokumente weltweit.

Die so erstellten Indexe werden dann im \www zur Verf"ugung gestellt. So kann man Suchanfragen der Form "`film and not movie"'  starten, um Dokumente die das Wort "`film"', nicht jedoch das Wort "`movie"' enthalten, als Ergebnis zu bekommen.


\section{Unzul"anglichkeiten von Suchmaschinen}
Das \www ist mittlerweile so gro"s, da"s Verbesserungen am Dokumentstandard HTML\footnote{http://www.w3.org/pub/WWW/MarkUp/} nur sehr tr"age Verbreitung finden. Die Ver"anderungen m"ussen zum einen in einem langwierigen Proze"s vom W3C verabschiedet werden, zum anderen mu"s daran anschlie"send die gesamte Software an die neuen Standards angepa"st werden.

Firmen wie Microsoft und Netscape, die den Web-Browser-Markt beherrschen, versuchen "uberdies noch eigene Standards mit Ihren Web-Browsern zu etablieren.

Aus diesem Grunde ist es sicher effizienter, Algorithmen f"ur eine bessere Aquirierung und Aufbereitung der Informationen, welche Suchmaschinen anbieten, zu verbessern.

Um die Suchmaschinen von ihren Unzul"anglichkeiten zu befreien, mu"s man sich derer erst einmal bewu"st werden.

Die heute eingesetzten Suchmaschinen ( z.B. Harvest\footnote{http://harvest.transarc.com/}, MomSpider\footnote{http://www.ics.uci.edu/pub/websoft/MOMspider/} oder Scooter\footnote{Scooter ist der von AltaVista, Lycos, Yahoo\dots \ eingesetzte, nicht frei verf"ubare Crawler.} ) laufen jeweils nur auf einem einzigen Rechner und versuchen das Web von diesem Rechner aus zu indexieren. Wenn man sich vorstellt, 150 Gigabyte an Text mit einem einzigen, wenngleich gro"sem Rechner aus dem Netz zu laden, kann man sich ausmalen, da"s diese Suchmaschinen bald an ihre Grenzen sto"sen.

Indizes kamen erstmals in B"uchern vor und sind sicher eine gro"se Hilfe um in B"uchern die gew"unschten Informationen schneller aufzufinden. Dies liegt zum einen an dem Hintergrundwissen "uber das Thema des Buches, das der Verlag einflie"sen l"a"st, zum anderen an dem Wissen von unrelevanten Informationen.

So wird in einem Buch das Wort "`das"' nur im Index zu finden sein, wenn es um den Artikel als solchen geht. In den meisten B"uchern werden f"ur das Themengebiet irrelevante Begriffe nicht indexiert.

Man kann nun nat"urlich dazu "ubergehen das \www mittels geschultem Personal indexieren zu lassen. In diese Richtung arbeiten auch schon einige Suchindizes, wie z.B. Lycos\footnote{http://www.lycos.de/} oder WebDe\footnote{http://www.web.de/}, die versuchen, den in Deutschland befindlichen Teil des Internet manuell zu indexieren.

Dies macht sicher insofern Sinn, als da"s Internetbenutzer in Deutschland zum Gro"steil die deutsche Sprache sprechen und wahrscheinlich nicht an Kochrezepten interessiert sind, welche in chinesisch verfa"st sind.

Mit der Gr"o"se des Internet und der Anzahl der dort zur Verf"ugung gestellten Dokumente, wird jedoch klar, da"s man eine manuelle und akutelle Indexierung nur f"ur ausgew"ahlte Informationsgebiete machen kann. Ein vollst"andiger Index des \www wird, allein durch die Fluktuation der Daten, nur mit maschineller Hilfe m"oglich und verwaltbar.


\section{Verbesserungsm"oglichkeiten}
Wie schon erw"ahnt, mu"s man versuchen die Strategien, die eine Suchmaschine ausmachen, zu verbessern. Sind die "`brute force"' Algorithmen der momentan verf"ugbaren Suchmaschinen f"ur kleine Datenmengen noch ausreichend, sto"sen diese Algorithmen mit der aktuellen Gr"o"se des Internet an ihre Grenzen.

Verbesserungsm"oglichkeiten fallen einem bei den genannten Unzul"anglichkeiten sofort ins Auge:
\begin{itemize}
	\item	Verteilung der Suchmaschinen auf mehrere Rechner.
	\item	Sprachenfilterung bei der Indexierung.
	\item 	Kontextabh"angige Indexierung.
\end{itemize}
Zu der Kontextabh"angigen Indexierung ist zu bemerken, da"s hier das Internet sicher die gr"o"ste M"oglichkeit einer elektronischen Datenverarbeitung bisher vers"aumt:

Geht man heute in eine Bibliothek, so kann man z.B. nach "`Mozart"' suchen und bekommt einige B"ucher "uber das Leben und Werk von Mozart. Durch geschickte elektronische Datenverarbeitung kann man jedoch den Informationssuchenden im Internet zum einen auf die Dokumente in denen "`Mozart"' behandelt wird verweisen, zum anderen kann man auch die mit "`Mozart"' verbundenen Themen aufzeigen. Der Informationssuchende kann dann gleich einen "Uberblick "uber das Themengebiet bekommen, ohne vorher mehrere B"ucher durchzulesen.


\section{Eingrenzung der Arbeit}
Im Rahmen dieser Diplomarbeit werden L"osungen f"ur die genannten Verbesserungsm"oglichkeiten aufgezeigt und implementiert.

Zu Beginn der Arbeit stand die Sprachenerkennung im Vordergrund. Um einen Sprachenerkenner implementieren zu k"onnen, ben"otigt man zun"achst einen Web-Crawler\footnote{to crawl: {\em engl.} kriechen. Web-Crawler bewegen sich, den Links, die sie auf Dokumenten vorfinden, folgend ( kriechend ) durchs \www . Genauer wird auf Crawler im Kapitel \ref{Crawler} eingegangen.}. Ein Web-Crawler ist ein Programm welches das \www durchsucht.

Aus diesem Grunde widme ich mich in den nun folgenden Kapiteln zun"achst den Crawlern, um dann, darauf aufbauend, die Sprachenerkennung und folgend die Kontexterkennung n"aher zu untersuchen.





\section{Grundbegriffe}

Hier nicht erw"ahnte Begriffe befinden sich im Glossar.

\subsection{Internet}
Eigentlich versteht man unter dem Begriff "`Internet"' nur das Netzwerk, welches mit dem {\em Internetprotokoll} arbeitet. Das {\em Internetprotokoll} wurde Mitte der 70'er Jahre vom amerikanischen Verteidigungsministerium in Auftrag gegeben; es sollte die Kommunikation von Milit"arcomputern nach einem atomaren Angriff sicherstellen.

Bis Ende der 80'er Jahre wurde das Internetprotokoll nur von den amerikanischen Milit"arbeh"orden und Universit"aten weltweit benutzt. Erst mit dem \www und dessen Verbreitung in den letzten Jahren, wurde das Internet zu {\em dem} gr"o"sten Computernetzwerk der Welt.


\subsection{URL}
Eine URL ( Uniform Resource Locator ) ist, wie der Name schon sagt, eine eindeutige Beschreibung einer Datenquelle. "Uber eine URL wird ein Dokument weltweit {\em eindeutig} benannt. Diese Eindeutigkeit der Namen erreicht man durch die drei verschiedenen Teile aus denen eine URL zusammengesetzt ist:

\begin{description}
	\item[protocol:]	Hier wird das "Ubertragungsprotokoll genannt, "uber welches die Daten zur Verf"ugung gestellt werden.
	\item[host:]	Dies ist der Name des Rechners, auf welchem die Daten liegen.
	\item[path:]	Dies ist der Pfadname, unter dem das Dokument auf dem Rechner gespeichert ist.
\end{description}

Eine URL mit Namen "`http://www.cis.uni-muenchen.de/index.html"' beschreibt also ein Dokument, welches auf dem  Rechner mit dem Namen "`www.cis.uni-muenchen.de"' unter dem Pfad "`index.html"' gespeichert ist und man via "`http"' abrufen kann.


\subsection{HTTP}
Das HTTP ( HyperText Transfer Protokol ) ist ein Protokoll, welches die Kommunikation von Web-Browsern mit Web-Servern beschreibt. Im wesentlichen besteht es aus einem Frage-Antwort Schema, welches ungef"ahr so aussieht:

Auf eine Anfrage des Clients: "`Gib mir das Hypertextdokument mit der folgenden URL\dots mit Hilfe des genannten Protokolls."'
\begin{verbatim}
GET http://www.cis.uni-muenchen.de/index.html HTTP/1.1
\end{verbatim}
Anwortet der Server: "'Hallo Client, hier kommt das Dokument \dots"'
\begin{verbatim}
HTTP/1.0 200 Document follows
MIME-Version: 1.0
Server: CERN/3.0
Date: Tuesday, 20-May-97 09:22:13 GMT
Content-Type: text/html
Content-Length: 3768
Last-Modified: Monday, 10-Mar-97 15:53:29 GMT


...
\end{verbatim}


\subsection{HTML}
Im \www werden Hypertextdokumente in HTML ( HyperText Markup Language ) publiziert. HTML ist eine Anwendung von SGML ( Standard General Markup Language ).

SGML entstand durch das Bestreben, formatierte Texte zwischen verschiedenen Rechnerplattformen austauschen zu k"onnen. Es handelt sich bei SGML um ein deskriptives Markup, d.h. es werden Textteile des Dokuments mit beschreibenden Marken versehen ( z.B. Titel, Author, "Uberschrift,\dots ).



\subsection{HTML-Dokument}
Ein HTML-Dokument, auch Webseite oder Hypertextdokument genannt, enth"alt neben dem Text noch weitere Informationen wie z.B. aktive Verweise, Textformatierung, Musik, Film, etc. .

Wie der Name schon sagt ist der Text in HTML verfa"st.


\subsection{\www}
\www war der Name des ersten Web-Browsers von Tim Berners-Lee. Heute ist der Begriff \www ein Synonym f"ur alle Hypertextdokumente die im Internet zur Verf"ugung gestellt werden.

Das \www-Programm von Tim Berners-Lee lief auf dem wenig verbreiteten NeXTSTEP Betriebssystem. Durchschlagenden Erfolg hatte das \www erst, als der Web-Browser "`Mosaic"' von NSCA auf Macintosh, Unix und Windows Systemen frei verf"ugbar wurde. 


\subsection{Web-Browser}
Ein Web-Browser stellt die in HTML geschriebenen Dokumente formatiert auf dem Bildschirm des Benutzers dar. Wenn man einen aktiven Verweis in einem Hypertextdokument mit der Maus anklickt, wird daraufhin das Dokument von der angegebenen Datenquelle geholt und dargestellt.



\subsection{Suchindex}
Suchindizes bestehen aus drei Teilen:
\begin{description}
\item[Crawler:]		Crawler wird das Programm genannt, welches das Hypertextdokumente aus dem Internet holt. Mehr dazu im Kapitel \ref{Crawler}.
\item[Indexierer:]	Dieses Program indexiert alle Seiten, welche der Crawler aus dem \www holt. Man mu"s die Daten aus zwei wesentlichen G"unden indexieren:
	\begin{description}
		\item[Datenkompression:] 	Ein Index ben"otigt meist weniger als $1\over 3$ der Gr"o"se des zugrundeliegenden Datenbestandes.
		\item[Geschwindigkeit:]		Ein Index mu"s nicht erst die gesamten Daten durchsuchen, sondern verwaltet alle in den Hyptertextdokumenten vorkommenden Daten in speziellen Datenstrukturen. Diese Datenstrukturen erm"oglichen es sofort, die Stelle zu finden, an der die gesuchten Informationen stehen.
		
		Die Zeit zum Auf\/finden der Informationen ist um Gr"o"senordnungen kleiner, als die zum linearen Durchsuchen der Daten.
	\end{description}
\item[SearchInterface:]	Damit Benutzer des \www die Informationen abrufen k"onnen, ben"otigt man noch eine Schnittstelle. Dies Schnittstelle nimmt zum einen die Suchanfragen des Benutzers an, zum anderen liefert es die Antworten des Indexierers in lesbarer Form zur"uck.
\end{description}